\documentclass[12pt]{article}
\usepackage[utf8]{inputenc}
\usepackage[T1]{fontenc}
\usepackage{hyperref}
\title{Foi Language Resources for K-12 Curriculum Development}
\author{Compiled by Claude Sonnet 5.5 (claude-sonnet-5-5) for Andr\'as Kornai}
\date{}
\begin{document}
\maketitle

\section{Introduction}
Foi (Foe, Mubi River; ISO 639-3 \texttt{foi}) is an East Kutubuan language spoken along Lake Kutubu and the Mubi River in the Southern Highlands Province of Papua New Guinea. Speaker estimates vary widely (about 6,000--10,000 depending on source). The search found few resources.

\section{Primary Educational Resources}
No primers, textbooks or school programmes were found in the search results.

\section{Linguistic and Reference Materials}
\begin{itemize}
\item Murray Rule, \emph{Foe} culture and language book including a grammar sketch and dictionary (224 pages, Chevron Niugini, 1993), as catalogued by Glottolog: \url{https://glottolog.org/resource/reference/id/50152}
\item Wikipedia, Foi language: \url{https://en.wikipedia.org/wiki/Foi_language}
\item Wikipedia, Kutubuan languages: \url{https://en.wikipedia.org/wiki/Kutubuan_languages}
\end{itemize}

\section{Religious and Community Resources}
Joshua Project reports Bible portions (1961--1968) and a New Testament (1978), with no complete Bible; the audio Bible and Jesus film are listed as unavailable.
\begin{itemize}
\item \url{https://joshuaproject.net/languages/foi}
\item \url{https://joshuaproject.net/people_groups/11220}
\end{itemize}

\section{Digital and Audio Resources}
None found in the search results.

\section{Curriculum Development Context}
Joshua Project notes limited access to education in the remote terrain and names literacy support as a need. The Rule dictionary and grammar sketch and the New Testament are the main known written bases for any native-language curriculum; their availability and current orthography were not verified.
\end{document}